\documentclass[11pt]{article}
\usepackage[margin=1in]{geometry}
\usepackage{amsmath,amssymb,amsthm}
\usepackage{booktabs}
\usepackage{longtable}
\usepackage{graphicx}
\usepackage{hyperref}
\usepackage{enumitem}
\usepackage{caption}
\captionsetup{font=small}

\title{CC-SPR Geometry-Family Empirical Study: Supplementary Material\\[6pt]
\large Extended Tables, Ablation Details, Diagnostics, and Reproducibility}
\author{}
\date{}

\begin{document}
\maketitle

\tableofcontents
\newpage

% ============================================================
\section{Repository and Data Inventory}
\label{sec:inventory}
% ============================================================

Repository root: \texttt{/cluster/VAST/kazict-lab/e/lesion\_phes/code/ccspr}.

\subsection{Geometry-family artifacts}
\begin{itemize}[leftmargin=1.5em]
\item \textbf{Results JSON}: \texttt{results/geometry\_family/geometry\_family\_results.json}
\item \textbf{Checkpoint}: \texttt{results/geometry\_family/checkpoint.json}
\item \textbf{Benchmark CSV}: \texttt{results/geometry\_family/tables/geometry\_family\_benchmark.csv}
\item \textbf{Diagnostics CSV}: \texttt{results/geometry\_family/tables/geometry\_family\_diagnostics.csv}
\item \textbf{Ablation CSV}: \texttt{results/geometry\_family/tables/geometry\_family\_ablation.csv}
\item \textbf{Timing CSV}: \texttt{results/geometry\_family/tables/geometry\_family\_timing.csv}
\end{itemize}

\subsection{Prior data artifacts (inherited)}
\begin{itemize}[leftmargin=1.5em]
\item \textbf{LUAD anchor}: \texttt{results/metrics.csv}, \texttt{results/main\_results.csv}
\item \textbf{LUAD multiseed}: \texttt{results/metrics\_luad\_multiseed.csv}, \texttt{results/paired\_stats\_luad\_multiseed.csv}
\item \textbf{GSE161711}: \texttt{data/cll\_venetoclax/processed\_matrix.tsv}
\item \textbf{GSE165087}: \texttt{data/cll\_rs\_scrna/data.h5ad}
\end{itemize}

\subsection{Source code for new backends}
All five geometry backends are implemented in a single dispatcher:
\begin{itemize}[leftmargin=1.5em]
\item \texttt{src/ccspr/geometry/distance.py} --- \texttt{build\_distance()} with modes: \texttt{euclidean}, \texttt{ricci}, \texttt{diffusion}, \texttt{phate\_like}, \texttt{dtm}
\item Shared infrastructure: \texttt{\_build\_knn\_affinity()} provides the $k$NN affinity matrix used by diffusion, PHATE-like, and DTM backends
\end{itemize}

% ============================================================
\section{Experimental Design}
\label{sec:design}
% ============================================================

\subsection{Dataset configurations}

\begin{table}[h]
\centering
\small
\begin{tabular}{llccccc}
\toprule
Dataset & Tissue & $n$ & Splits & Bootstraps & Ablation & Bounded \\
\midrule
LUAD & TCGA lung adenocarcinoma & 60 & 2 & 5 & $\lambda$ sweep & No \\
GSE161711 & CLL bulk RNA-seq & 96 & 5 (CV) & 5 & $\lambda$ sweep & No \\
GSE165087 & CLL scRNA-seq & 320 & 1 & 1 & $\lambda$ sweep & Yes \\
\bottomrule
\end{tabular}
\caption{Dataset configurations for the geometry-family study.
LUAD uses \texttt{max\_samples=60} for tractability; GSE165087 is bounded to 320 cells.}
\label{tab:configs}
\end{table}

\subsection{Backends tested}

All five backends share the same downstream pipeline: Vietoris--Rips filtration $\to$ persistent homology ($H_1$) $\to$ harmonic/CC-SPR representative $\to$ TIP projection $\to$ random forest classification.
The only difference is the pairwise distance matrix fed to the Rips complex:

\begin{enumerate}[leftmargin=1.5em]
\item \textbf{Euclidean}: $d_{ij} = \|x_i - x_j\|_2$
\item \textbf{Ollivier--Ricci}: Curvature-adjusted Euclidean distances via discrete Ricci flow
\item \textbf{Diffusion}: Truncated eigendecomposition of the Markov transition matrix from a $k$NN affinity graph; $L^2$ distance in eigenspace
\item \textbf{PHATE-like}: Potential distance $\|-\log(P^t + \varepsilon)\|_2$ where $P$ is the Markov matrix
\item \textbf{DTM}: Distance-to-measure weighted: $\sqrt{d_{ij}^2 + \text{dtm}_i^2 + \text{dtm}_j^2}$, where $\text{dtm}_i$ is the mean distance to $k$ nearest neighbors
\end{enumerate}

% ============================================================
\section{Full Benchmark Results}
\label{sec:benchmark}
% ============================================================

\begin{table}[h]
\centering
\small
\begin{tabular}{llccl}
\toprule
Dataset & Backend & Mean F1 & Std F1 & Elapsed (s) \\
\midrule
LUAD & euclidean & 0.4342 & 0.0394 & 0.5 \\
LUAD & ricci & 0.5012 & 0.1073 & 5.2 \\
LUAD & diffusion & 0.5000 & 0.1000 & 0.5 \\
LUAD & phate\_like & 0.3546 & 0.0880 & 0.2 \\
LUAD & dtm & 0.5315 & 0.0765 & 0.2 \\
\midrule
GSE161711 & euclidean & 0.7483 & 0.0825 & 8.0 \\
GSE161711 & ricci & 0.5944 & 0.0845 & 29.0 \\
GSE161711 & diffusion & 0.5553 & 0.0721 & 2.5 \\
GSE161711 & phate\_like & 0.5396 & 0.0632 & 1.2 \\
GSE161711 & dtm & 0.6178 & 0.1333 & 14.4 \\
\midrule
GSE165087 & euclidean & 0.9254 & 0.0 & 1099.0 \\
GSE165087 & ricci & 0.9254 & 0.0 & 6.2 \\
GSE165087 & diffusion & 0.9254 & 0.0 & 1.9 \\
GSE165087 & phate\_like & 0.9254 & 0.0 & 36.3 \\
GSE165087 & dtm & 0.9254 & 0.0 & 3932.1 \\
\bottomrule
\end{tabular}
\caption{Complete geometry-family benchmark results.
All 15 dataset$\times$backend combinations completed successfully via Slurm job 12883024.
GSE165087 achieves identical F1 across all backends due to bounded sample size (320 cells, 1 split).}
\label{tab:benchmark_full}
\end{table}

% ============================================================
\section{Full Diagnostics}
\label{sec:diagnostics}
% ============================================================

\begin{table}[h]
\centering
\small
\begin{tabular}{llccccccc}
\toprule
Dataset & Backend & $n_\text{ok}$ & Lifetime & $\sigma_\text{life}$ & Prominence & $\sigma_\text{prom}$ & Support & Entropy \\
\midrule
LUAD & euclidean & 5 & 0.0798 & 0.025 & 0.0469 & 0.0296 & 50 & 3.733 \\
LUAD & ricci & 4 & 0.0130 & 0.007 & 0.0067 & 0.008 & 71 & 4.226 \\
LUAD & diffusion & 5 & 0.0256 & 0.018 & 0.0146 & 0.015 & 78 & 4.274 \\
LUAD & phate\_like & 4 & 0.0112 & 0.010 & 0.0085 & 0.008 & 80 & 4.382 \\
LUAD & dtm & 5 & 0.0083 & 0.006 & 0.0068 & 0.007 & 76 & 4.273 \\
\midrule
GSE161711 & euclidean & 5 & 0.0828 & 0.003 & 0.0118 & 0.008 & 59 & 4.011 \\
GSE161711 & ricci & 5 & 0.1248 & 0.048 & 0.0468 & 0.041 & 95 & 4.536 \\
GSE161711 & diffusion & 5 & 0.0426 & 0.024 & 0.0213 & 0.021 & 93 & 4.508 \\
GSE161711 & phate\_like & 5 & 0.0091 & 0.007 & 0.0049 & 0.005 & 84 & 4.373 \\
GSE161711 & dtm & 5 & 0.0487 & 0.008 & 0.0166 & 0.006 & 24 & 3.130 \\
\bottomrule
\end{tabular}
\caption{Full topological diagnostics for LUAD and GSE161711.
$n_\text{ok}$ is the number of bootstrap runs producing valid $H_1$ bars.
Support is the number of nonzero TIP entries.
GSE165087 diagnostics omitted (single bootstrap, bounded configuration).}
\label{tab:diagnostics_full}
\end{table}

\subsection{Diagnostic interpretation}

\paragraph{Lifetime and prominence.}
Euclidean consistently produces the longest-lived $H_1$ bars (LUAD: 0.080, GSE161711: 0.083), suggesting that the Euclidean metric preserves larger-scale topological features.
Curvature-based backends (Ricci, diffusion) produce shorter-lived but more numerous cycles, consistent with their tendency to redistribute edge lengths and create more uniform filtration scales.

\paragraph{Support size.}
PHATE-like achieves the broadest TIP support on LUAD (80 features) but low classification F1 (0.355), suggesting that distributing representative weight across many features dilutes discriminative signal.
Conversely, Euclidean has the smallest support (50) on LUAD, concentrating signal on fewer genes.

\paragraph{TIP entropy.}
Higher entropy indicates more uniform distribution of topological influence across features.
The curvature-based backends consistently achieve higher entropy (4.2--4.5) than Euclidean (3.7--4.0), reflecting their role in ``flattening'' the metric landscape.

\paragraph{TIP Gini coefficient.}
\begin{table}[h]
\centering
\small
\begin{tabular}{lcc}
\toprule
Backend & LUAD Gini & GSE161711 Gini \\
\midrule
euclidean & 0.321 & 0.186 \\
ricci & 0.098 & 0.047 \\
diffusion & 0.180 & 0.065 \\
phate\_like & 0.000 & 0.136 \\
dtm & 0.164 & 0.130 \\
\bottomrule
\end{tabular}
\caption{TIP Gini coefficients by backend.
Lower Gini indicates more uniform TIP distribution.
Ricci achieves the most uniform distribution on both datasets; PHATE-like achieves perfect uniformity on LUAD (Gini = 0).}
\label{tab:gini}
\end{table}

% ============================================================
\section{Lambda Ablation Details}
\label{sec:ablation}
% ============================================================

The geometry-family ablation sweeps the sparsity parameter $\lambda \in \{10^{-6}, 10^{-5}\}$ across all five backends and three datasets.

\begin{table}[h]
\centering
\small
\begin{tabular}{llccc}
\toprule
Dataset & Backend & $\lambda = 10^{-6}$ F1 & $\lambda = 10^{-5}$ F1 & $\Delta$ \\
\midrule
LUAD & euclidean & 0.423 & 0.423 & 0.000 \\
LUAD & ricci & 0.463 & 0.463 & 0.000 \\
LUAD & diffusion & 0.380 & 0.380 & 0.000 \\
LUAD & phate\_like & 0.496 & 0.496 & 0.000 \\
LUAD & dtm & 0.530 & 0.530 & 0.000 \\
\midrule
GSE161711 & euclidean & 0.646 & 0.595 & $-0.051$ \\
GSE161711 & ricci & 0.542 & 0.522 & $-0.021$ \\
GSE161711 & diffusion & 0.537 & 0.508 & $-0.029$ \\
GSE161711 & phate\_like & 0.696 & 0.696 & 0.000 \\
GSE161711 & dtm & 0.588 & 0.628 & $+0.040$ \\
\midrule
GSE165087 & all backends & 0.913 & 0.913 & 0.000 \\
\bottomrule
\end{tabular}
\caption{Lambda ablation results.
LUAD is insensitive to the tested $\lambda$ range.
GSE161711 shows mild sensitivity: Euclidean and diffusion prefer lower $\lambda$; DTM prefers higher $\lambda$.
GSE165087 is invariant (bounded, single split).}
\label{tab:ablation_full}
\end{table}

\subsection{Ablation interpretation}
The narrow $\lambda$ range ($10^{-6}$ to $10^{-5}$) was chosen to probe sensitivity near the default operating point.
LUAD shows zero sensitivity across all backends, suggesting that the sparse representative solution is stable in this regime.
GSE161711 shows mild but non-trivial sensitivity for Euclidean ($-0.051$) and DTM ($+0.040$), indicating that the optimal sparsity level is backend-dependent even within a single dataset.

% ============================================================
\section{Timing Analysis}
\label{sec:timing}
% ============================================================

\begin{table}[h]
\centering
\small
\begin{tabular}{llr}
\toprule
Dataset & Backend & Elapsed (s) \\
\midrule
LUAD (60 samples) & euclidean & 0.5 \\
 & ricci & 5.2 \\
 & diffusion & 0.5 \\
 & phate\_like & 0.2 \\
 & dtm & 0.2 \\
\midrule
GSE161711 (96 samples) & euclidean & 8.0 \\
 & ricci & 29.0 \\
 & diffusion & 2.5 \\
 & phate\_like & 1.2 \\
 & dtm & 14.4 \\
\midrule
GSE165087 (320 samples) & euclidean & 1099.0 \\
 & ricci & 6.2 \\
 & diffusion & 1.9 \\
 & phate\_like & 36.3 \\
 & dtm & 3932.1 \\
\bottomrule
\end{tabular}
\caption{Wall-clock timing for each dataset$\times$backend combination.
DTM on GSE165087 (3932s $\approx$ 65 min) is the most expensive combination due to the DTM-weighted distance creating a dense filtration on 320 samples.
Euclidean on GSE165087 (1099s $\approx$ 18 min) is expensive due to direct pairwise distance computation at scale.}
\label{tab:timing_full}
\end{table}

\subsection{Scaling observations}
\begin{itemize}[leftmargin=1.5em]
\item Diffusion and PHATE-like backends are consistently fast ($< 40$s even at 320 samples) because they operate in a reduced eigenspace.
\item Ricci is moderate ($5$--$29$s) but does not scale quadratically because the curvature computation uses $k$NN sparsity.
\item Euclidean scales quadratically with $n$ (0.5s at $n=60$ to 1099s at $n=320$).
\item DTM inherits Euclidean scaling and adds overhead from neighbor averaging, making it the slowest backend at large $n$.
\end{itemize}

% ============================================================
\section{Execution Envelope}
\label{sec:execution}
% ============================================================

\subsection{Geometry-family Slurm job}

\begin{table}[h]
\centering
\small
\begin{tabular}{lllll}
\toprule
Job ID & Name & State & Walltime & Resources \\
\midrule
12883024 & ccspr\_geometry\_family & COMPLETED & $< 12$h & 48\,GB, 8 CPUs \\
\bottomrule
\end{tabular}
\caption{Slurm execution for the geometry-family experiment.
All 15 backend$\times$dataset combinations completed, plus diagnostics and ablation sweeps.}
\label{tab:geom_job}
\end{table}

\subsection{Prior Slurm history (inherited)}
See prior supplement for the complete job history including the scRNA OOM/timeout failure sequence and manuscript-safe resolution.

% ============================================================
\section{Backend Implementation Details}
\label{sec:impl}
% ============================================================

\subsection{Shared kNN affinity}
The diffusion, PHATE-like, and DTM backends share a common $k$NN affinity construction:
\begin{enumerate}[leftmargin=1.5em]
\item Compute $k$ nearest neighbors ($k=15$ by default) using \texttt{sklearn.neighbors.NearestNeighbors}
\item Build adaptive Gaussian kernel: $W_{ij} = \exp(-d_{ij}^2 / \sigma_i \sigma_j)$ where $\sigma_i$ is the distance to the $k$-th neighbor of point $i$
\item Symmetrize: $W \leftarrow (W + W^\top) / 2$
\end{enumerate}

\subsection{Diffusion distance}
From the affinity matrix $W$:
\begin{enumerate}[leftmargin=1.5em]
\item Form row-stochastic Markov matrix $P = D^{-1}W$ where $D = \text{diag}(W\mathbf{1})$
\item Compute top $m=50$ eigenvectors of $P$ via \texttt{scipy.sparse.linalg.eigsh}
\item Diffusion coordinates: $\Psi_i = [\lambda_1 \psi_1(i), \ldots, \lambda_m \psi_m(i)]$
\item Distance: $d_\text{diff}(i,j) = \|\Psi_i - \Psi_j\|_2$
\end{enumerate}

\subsection{PHATE-like potential distance}
\begin{enumerate}[leftmargin=1.5em]
\item Compute Markov matrix $P$ as above
\item Raise to power $t=5$: $P^t$ (diffusion time)
\item Potential: $U = -\log(P^t + \varepsilon)$ with $\varepsilon = 10^{-7}$
\item Distance: $d_\text{phate}(i,j) = \|U_i - U_j\|_2$
\end{enumerate}
This is a lightweight surrogate for full PHATE; it captures the same potential-distance intuition without the full PHATE embedding pipeline.

\subsection{DTM-weighted distance}
\begin{enumerate}[leftmargin=1.5em]
\item For each point $i$, compute $\text{dtm}_i = \frac{1}{k}\sum_{j \in \text{kNN}(i)} d(i,j)$
\item Weighted distance: $d_\text{dtm}(i,j) = \sqrt{d_{ij}^2 + \text{dtm}_i^2 + \text{dtm}_j^2}$
\end{enumerate}
The DTM weight inflates distances near sparse regions (high local $\text{dtm}$) and preserves distances in dense regions, providing noise robustness.

% ============================================================
\section{GSE165087 Ceiling Effect}
\label{sec:ceiling}
% ============================================================

All five backends achieve identical F1 (0.9254) on GSE165087.
This ceiling effect has two causes:
\begin{enumerate}[leftmargin=1.5em]
\item \textbf{Bounded sample}: 320 cells with 1 train/test split and 1 bootstrap provide insufficient statistical resolution to differentiate backends.
\item \textbf{Strong class signal}: The cell-type labels in this CLL scRNA dataset are well-separated in PCA space, so the random forest classifier saturates regardless of which topological features are provided.
\end{enumerate}
This result is reported honestly as a feasibility check, not as evidence that backends are equivalent.

% ============================================================
\section{Cross-Dataset Pattern Interpretation}
\label{sec:patterns}
% ============================================================

\paragraph{Backend ranking is dataset-dependent.}
On LUAD, DTM $>$ Ricci $>$ Diffusion $>$ Euclidean $>$ PHATE-like.
On GSE161711, Euclidean $>$ DTM $>$ Ricci $>$ Diffusion $>$ PHATE-like.
The rank inversion between Euclidean and curvature-based backends is the central empirical finding.

\paragraph{Why curvature helps on LUAD.}
LUAD ($n=60$, high-dimensional TCGA expression) benefits from curvature correction because Euclidean neighborhood graphs contain shortcut edges between lung adenocarcinoma subtypes.
Ricci flow and DTM weighting remove these shortcuts by penalizing edges in sparse inter-cluster regions.

\paragraph{Why Euclidean dominates on GSE161711.}
GSE161711 ($n=96$, lower-dimensional bulk RNA-seq) has cleaner class separation in the original metric space.
Curvature-based transformations introduce noise at this sample size because the $k$NN graph is less stable.

\paragraph{Support--accuracy tradeoff.}
Backends with higher TIP support (phate\_like: 80, diffusion: 78 on LUAD) tend to have lower F1, while backends with concentrated support (euclidean: 50) or curvature-sharpened support (ricci: 71, dtm: 76) achieve better classification.
This suggests a ``dilution vs.\ concentration'' tradeoff in topological feature projection.

% ============================================================
\section{Reproducibility}
\label{sec:repro}
% ============================================================

\subsection{Environment}
\begin{verbatim}
source /cluster/VAST/kazict-lab/e/lesion_phes/lesenv/bin/activate
export PYTHONPATH=src
\end{verbatim}

\subsection{Geometry-family experiment}
\begin{verbatim}
# Submit to Slurm
sbatch slurm/hellbender_geometry_family.sbatch

# Or run directly (requires 48GB+ RAM)
python3.11 -u experiments/run_geometry_family.py
\end{verbatim}

\subsection{Checkpointing}
The experiment runner saves a checkpoint after each phase (LUAD, GSE161711, GSE165087).
If interrupted, re-running the script will skip completed phases by reading \texttt{results/geometry\_family/checkpoint.json}.

\end{document}
